% p.71 — العوامل ومتعلقاتها (تتمة)
% Layout: twocol. rightcol: entries 80–83, rule (افعال المدح والذم), 84–87, rule (افعال القلوب),
% 88–94, rule (منها الفعل …); leftcol: 95–101, rule (الاسم المبهم), 102, rule (الابتداء), 103–104,
% the small note on التوابع, 105–106 (106 continues on p.72). Rule paragraphs are \par-separated
% text between the sarf tables.
% Note: reledpar sets the two columns line by line (line i of the right column beside line i of
% the left one, each pair as tall as its taller member), so a column made of tall sarf "lines"
% padded the other column and the page overflowed. Each column's whole content is therefore
% wrapped in one rightbox/leftbox (minipages inside the reledpar columns): the columns then
% flow independently and the page fits; the apparatus numbers each column as a single line.
\mspage{71}
\begin{twocol}
\begin{rightbox}
\begin{sarf}{|c|c|}
(80) عَسٰى & عَسٰى زَيْدٌ اَنْ يَخْرُجَ\\
 & عَسَتْ هِنْدَةُ اَنْ تَخْرُجَ\\
(81) كَادَ & كَادَ زَيْدٌ يَخْرُجُ\\
(82) كَرَبَ & كَرَبَ زَيْدٌ يَخْرُجُ\\
(83) اَوْشَكَ & اَوْشَكَ زَيْدٌ يَخْرُجُ\\
 & اَوْشَكَ زَيْدٌ اَنْ يَخْرُجَ\\
\end{sarf}
مِنْهَا اَفْعَالُ الْمَدْحِ وَالذَّمِّ تَرْفَعُ الْفَاعِلَ الْمُعَرَّفَ بِالْاَلِفِ وَاللَّامِ وَهِيَ اَرْبَعَةٌ
\par
\begin{sarf}{|c|c|}
(84) نِعْمَ & نِعْمَ الرَّجُلُ زَيْدٌ نِعْمَ رَجُلًا زَيْدٌ\\
 & نِعْمَ غُلَامُ الرَّجُلِ زَيْدٌ نِعْمَ الرَّجُلُ\\
(85) بِئْسَ & بِئْسَ الرَّجُلُ زَيْدٌ\\
(86) سَاءَ & سَاءَ الرَّجُلُ زَيْدٌ\\
(87) حَبَّذَا & حَبَّذَا زَيْدٌ حَبَّذَا فَاطِمَةُ\\
\end{sarf}
مِنْهَا اَفْعَالُ الشَّكِّ وَالْيَقِيْنِ تُسَمّٰى اَفْعَالَ الْقُلُوْبِ تَرْفَعُ الْفَاعِلَ وَتَنْصِبُ مَفْعُوْلَيْنِ وَهِيَ سَبْعَةٌ
\par
\begin{sarf}{|c|c|}
(88) حَسِبْتُ & حَسِبَ زَيْدٌ عَمْرًا عَالِمًا\\
(89) خِلْتُ & خَالَ زَيْدٌ خَالِدًا قَائِمًا\\
(90) ظَنَنْتُ & ظَنَّ زَيْدٌ بَكْرًا نَائِمًا\\
(91) رَاَيْتُ & رَاَيْتُ زَيْدًا فَقِيْهًا فَانْظُرْ مَاذَا تَرٰى\\
(92) وَجَدْتُ & وَجَدَ زَيْدٌ الْجِدَارَ سَاقِطًا\\
 & وَجَدْتُ الضَّالَّةَ\\
(93) عَلِمْتُ & عَلِمْتُ زَيْدًا فَاضِلًا - اَعْلَمْتُ زَيْدًا عَمْرًا فَاضِلًا\\
(94) زَعَمْتُ & زَعَمَ عَبْدُ اللهِ زَيْدًا فَاسِقًا\\
\end{sarf}
مِنْهَا الْفِعْلُ تَرْفَعُ الْفَاعِلَ فَقَطْ اَوْ تَنْصِبُ الْمَفْعُوْلَ اَيْضًا
\par
\end{rightbox}
\begin{leftbox}
\begin{sarf}{|c|c|}
(95) اَلْفِعْلُ & ذَهَبَ زَيْدٌ - ضَرَبَ زَيْدٌ عَمْرًا\\
(96) اَلْفِعْلُ الْمَجْهُوْلُ & ضُرِبَ عَمْرٌو\\
(97) اَلْاِسْمُ الْفَاعِلُ & زَيْدٌ ذَاهِبٌ غُلَامُهُ (لَازِمٌ)\\
 & زَيْدٌ ضَارِبٌ غُلَامُهُ عَمْرًا (مُتَعَدٍّ)\\
(98) اَلْاِسْمُ الْمَفْعُوْلُ & زَيْدٌ مَضْرُوْبٌ غُلَامُهُ\\
(99) اَلصِّفَةُ الْمُشَبَّهَةُ & زَيْدٌ حَسَنٌ وَجْهُهُ\\
(100) اَلْمَصْدَرُ & عَجِبْتُ مِنْ ذَهَابِ زَيْدٍ\\
 & عَجِبْتُ مِنْ ضَرْبِ زَيْدٍ عَمْرًا (مَعْرُوْفٌ)\\
 & عَجِبْتُ مِنْ ضَرْبِ عَمْرٍو (مَجْهُوْل)\\
(101) اَلْمُضَافُ & جَاءَنِيْ غُلَامُ زَيْدٍ\\
\end{sarf}
مِنْهَا الْاِسْمُ الْمُبْهَمُ لَا يَظْهَرُ مَعْنَاهُ اِلَّا بِالتَّمْيِيْزِ
\par
\begin{sarf}{|c|c|}
(102) اَلْاِسْمُ الْمُبْهَمُ & عِنْدِيْ رَطْلٌ زَيْتًا\\
 & عِنْدِيْ مَنَوَانِ سَمَنًا\\
 & عِنْدِيْ عِشْرُوْنَ دِرْهَمًا\\
 & \unclear{اِنْ لِزَيْدٍ مَلَأَ دَلْوٍ عَسَلًا}\\ % TODO: first words smudged; perhaps لِزَيْدٍ مِلْءُ دَلْوٍ عَسَلًا
\end{sarf}
مِنْهَا الْاِبْتِدَاءُ يَرْفَعُ الْاِسْمَ وَالْخَبَرَ
\par
\begin{sarf}{|c|c|}
(103) اَلْاِبْتِدَاءُ & زَيْدٌ قَائِمٌ فِي الدَّارِ زَيْدٌ\\
(104) \unclear{خُلُوٌّ} & يَضْرِبُ لَمْ يَضْرِبْ لَنْ يَضْرِبَ\\
\end{sarf}
{\small اعلم ان التوابع خمسة اعرابها تابع للمتبوع}
\par
\begin{sarf}{|c|c|}
(105) صِفَةٌ & هٰذَا رَجُلٌ عَالِمٌ رَاَيْتُ رَجُلًا عَالِمًا\\
(106) بَدَلٌ & رَاَيْتُ زَيْدًا صَدِيْقَكَ مَرَرْتُ بِزَيْدٍ \unclear{سَيِّدِكَ}\\ % TODO: last word cut off at the margin; entry continues on p.72
\end{sarf}
\end{leftbox}
\end{twocol}
